\scanpage{60}
\uline{fonction elliptique}. C'est une fonction méromorphe dans le plan à distance finie, et admettant deux périodes \unsure{distinctes|distincte} $\omega_1$ et $\omega_2$. En particulier, une constante est une fonction elliptique. Si $f(z)$ n'est pas une constante, le rapport de $\omega_1$ et $\omega_2$ n'est \deleted{imaginaire} \inserted{pas réel}.

\inserted{\unsure{Définition} de la période ($\omega_1$ et $\omega_2$ \unsure{donnés|donnent} une fois pour \unsure{ttes}) La donnée d'une fonction elliptique n'entraîne \uline{pas} la donnée de $\omega_1$ et $\omega_2$. Mais il faut remarquer que}

Pavage correspondant du plan. On peut définir un réseau de parallélogrammes \unsure{correspondants|équipollents} aux périodes $\omega_1$ et $\omega_2$ et à un sommet $a$ arbitraire. Ils sont \unsure{déterminé|déterminés} par leurs sommets de la forme $a+n_1\omega_1+n_{\unsure{1|2}}\omega_2$ ; tout point qui n'est pas de la forme $a+n\omega_1+\lambda\omega_1$, ou $a+\lambda\omega_1+n\omega_2$ ($\unsure{b|\lambda}$ et $n$ réels, $n$ entier) appartient : 1 et 1 seul de ces parallélogrammes, \inserted{\unsure{il lui est intérieur|il lui est int.}} \unsure{le pt|tout point} qui est de la forme précédente, mais non de la forme $a+n_1\omega_1\unsure{-|+}n_1\omega_2=a+\omega$, \deleted{est} appartient à deux et deux seult de ces par. \deleted{\unsure{en}} \inserted{il appartient aux côtés, non au sommet} enfin, un point de la forme $a\unsure{-|+}n_1\omega_1\unsure{-|+}n_1\omega_2$ appartient à 4 et 4 seult de tels par, il appartient à leurs sommets.

Si on se donne un parallélogramme de périodes, tout point est l'``homologue'' soit d'un pt intérieur, soit de deux p\textsuperscript{ts} sur les côtés, et non les sommets, soit des 4 sommets. \inserted{En particulier, deux p\textsuperscript{ts} \uline{int} à un $\overline{\mathrm m}$ par. de périodes ne sont pas homologues} \deleted{Si on se donne la fonction}

On peut trouver des par. des périodes, tels que sur leur côté ne se trouve aucun zéro, pôle etc. ; Dans ce cas, la \deleted{\unsure{nb}} \inserted{liste} des zéros (pôles) avec leur ordre de multiplicité, est \unsure{indépendant|indépendante} du par. considéré. Il en est ainsi en tte généralité, si on \unsure{cpte|compte} \deleted{une seule fois} \deleted{\unsure{disting}}é les zéros de tout groupe de zéros homologues (\deleted{\unsure{situé|situés}} ils est situés sur les sommets - auquel cas ils sont 4, on les divise par 4 -- autrement on les divise par deux)

% archive sheet number: 59
